\documentclass[12pt,a4paper]{article}

\usepackage[margin=2.5cm]{geometry}
\usepackage{hyperref}
\usepackage{listings}
\usepackage{xcolor}
\usepackage{booktabs}
\usepackage{enumitem}
\usepackage{parskip}
\usepackage{titlesec}

\hypersetup{colorlinks=true, linkcolor=blue, urlcolor=blue}

\lstset{
  basicstyle=\ttfamily\small,
  keywordstyle=\color{blue},
  commentstyle=\color{gray},
  stringstyle=\color{orange!80!black},
  breaklines=true,
  frame=single,
  tabsize=2
}

\title{\textbf{Recipe Maker — Project Report}\\[0.4em]
  \large Patient Meal Ordering with Component Customisation and Search Strategies}
\author{Software Engineering — Academic Year 2025/2026}
\date{11 May 2026}

\begin{document}
\maketitle
\tableofcontents
\newpage

% ─────────────────────────────────────────────────────────────
\section{Project Overview}
Recipe Maker is a full-stack hospital meal-management application built with
\textbf{Spring Boot 3} (backend) and \textbf{React 18 + TypeScript} (frontend).
It allows dieticians to define meal recipes composed of a mandatory
\emph{main component} and zero or more optional \emph{modifiable components}.
Doctors assign medical conditions to patient profiles, and the system enforces
ingredient compatibility at ordering time so that a patient can never receive
a component that conflicts with their conditions.

\subsection{Technology Stack}
\begin{center}
\begin{tabular}{ll}
\toprule
\textbf{Layer} & \textbf{Technology} \\
\midrule
Backend      & Java 21, Spring Boot 3, Spring Security, JPA / Hibernate \\
Database (runtime) & PostgreSQL 16-alpine (Docker volume \texttt{pgdata}) \\
Database (tests)   & H2 in-memory, PostgreSQL compatibility mode \\
Frontend     & React 18, TypeScript 5, Vite 5.2, React Router v6 \\
HTTP client  & Axios 1.7 (Basic Auth interceptor) \\
WebSocket    & STOMP / SockJS (\texttt{@stomp/stompjs} + \texttt{sockjs-client}) \\
API style    & REST (JSON), Springdoc / OpenAPI, Swagger UI \\
Deployment   & Docker Compose (app service + PostgreSQL service) \\
\bottomrule
\end{tabular}
\end{center}

% ─────────────────────────────────────────────────────────────
\section{Domain Model}

\subsection{Core Entities}
\begin{itemize}[noitemsep]
  \item \texttt{Recipe} — a meal with a course (BREAKFAST / LUNCH / DINNER), a type
        (MAIN / SIDE), one \emph{non-modifiable} main component, an optional
        set of \emph{modifiable} components, and an \texttt{imageUrl} field
        (Unsplash CDN URL rendered as a card banner in the UI).
  \item \texttt{RecipeComponent} — abstract JPA entity with single-table inheritance.
        Concrete subtypes: \texttt{ModifiableComponent} and \texttt{NonModifiableComponent}.
        Each component carries a \texttt{ManyToMany} relationship to
        \texttt{PatientCondition} to express incompatibilities.
  \item \texttt{PatientProfile} — linked to an \texttt{AppUser} with role
        \texttt{PATIENT}; holds a set of \texttt{PatientCondition}s.
  \item \texttt{MealOrder} — records a patient's order for a recipe together with
        the specific modifiable components they chose.
\end{itemize}

\subsection{Component Incompatibility}
The join table \texttt{component\_incompatibility} associates a
\texttt{RecipeComponent} with one or more \texttt{PatientCondition}s.
At ordering time the service validates \textbf{two layers}:
\begin{enumerate}[noitemsep]
  \item The recipe's \emph{non-modifiable main component} must not conflict with any
        of the patient's conditions.  If it does, the order is immediately rejected
        with \texttt{400 Bad Request}.
  \item For each optional (modifiable) component the patient selected, none of its
        incompatible conditions may appear in the patient's profile.
\end{enumerate}

% ─────────────────────────────────────────────────────────────
\section{Design Patterns}

\subsection{Observer Pattern — Order Lifecycle Notifications}
\texttt{OrderEventObserver} is an interface with two methods:
\texttt{onOrderPlaced(MealOrder)} and
\texttt{onOrderStatusAdvanced(MealOrder, OrderStatus)}.  \texttt{MealOrderService}
maintains a list of observers and notifies them after each state change.
\texttt{NotificationService} is the sole registered observer; it persists a
\texttt{Notification} entity and pushes it to the patient via STOMP/WebSocket.
Adding further observers (e.g.\ email or SMS dispatch) requires no change to
\texttt{MealOrderService}.

\subsection{Strategy Pattern — Recipe Search}
Both the backend and the frontend implement the \emph{Strategy} pattern for
recipe filtering.

\subsubsection{Backend (\texttt{com.example.recipemaker.search})}
\begin{itemize}[noitemsep]
  \item \texttt{RecipeSearchStrategy} — interface with a single \texttt{matches(Recipe)} method.
  \item \texttt{NameSearchStrategy} — case-insensitive substring match on recipe name.
  \item \texttt{ComponentNameSearchStrategy} — matches if any component name contains the keyword.
  \item \texttt{ConditionCompatibilityStrategy} — filters out recipes whose main component
        conflicts with the patient's conditions.
  \item \texttt{CompositeAndSearchStrategy} — combines strategies with AND semantics.
  \item \texttt{CompositeOrSearchStrategy} — combines strategies with OR semantics.
\end{itemize}

\subsubsection{Frontend (\texttt{src/patterns/searchStrategy.ts})}
The frontend mirrors the backend hierarchy so that the recipe list can be
filtered client-side without extra round-trips:
\begin{itemize}[noitemsep]
  \item \texttt{RecipeSearchStrategy} (interface)
  \item \texttt{NameSearchStrategy}
  \item \texttt{ComponentNameSearchStrategy}
  \item \texttt{MealTypeSearchStrategy} — \emph{newly added}; filters by MAIN or SIDE.
  \item \texttt{CompositeAndSearchStrategy}
  \item \texttt{CompositeOrSearchStrategy}
\end{itemize}

The factory function \texttt{buildAndStrategy(filters)} accepts \texttt{name},
\texttt{componentName}, and the new \texttt{mealType} parameter, composes the
appropriate leaf strategies, and returns a \texttt{CompositeAndSearchStrategy}.

\subsection{Builder Pattern}
\texttt{MealOrder}, \texttt{MealOrderResponse}, and other DTOs use Lombok's
\texttt{@Builder} annotation.  The frontend mirrors this in
\texttt{src/patterns/recipeRequestBuilder.ts}.

\subsection{Repository Pattern}
API access on the frontend is encapsulated in dedicated repository objects
(\texttt{orderRepository}, \texttt{recipeRepository}, etc.) that expose typed
async methods.  On the backend, Spring Data JPA repositories provide the same
abstraction.

% ─────────────────────────────────────────────────────────────
\section{Feature: Real-Time Notifications (WebSocket / STOMP)}

\subsection{Motivation}
Kitchen staff advance order statuses via the \texttt{/api/kitchen/orders/\{id\}/advance}
endpoint.  Patients need to see status updates without polling.

\subsection{Backend}
\texttt{WebSocketConfig} registers a STOMP endpoint at \texttt{/ws} (SockJS
enabled) and configures a simple in-memory message broker.  After every
\texttt{onOrderStatusAdvanced} call, \texttt{NotificationService} uses
\texttt{SimpMessagingTemplate} to push a message to
\texttt{/user/\{username\}/queue/notifications}.

Authentication is handled via the STOMP CONNECT frame's \texttt{login} /
\texttt{passcode} headers, validated by a custom \texttt{ChannelInterceptor}.

\subsection{Frontend}
\texttt{NotificationContext.tsx} establishes the SockJS/STOMP connection on mount,
subscribes to \texttt{/user/queue/notifications}, and appends incoming messages
to a React state list.  \texttt{NotificationBell.tsx} shows the unread count and
a dropdown of recent notifications.

% ─────────────────────────────────────────────────────────────
\section{Feature: User Self-Registration}

\subsection{Motivation}
Previously there was no UI path for new users to register without
direct database access or Swagger UI.  The updated system exposes a
dedicated registration dialog on the login page.

\subsection{Backend}
\texttt{POST /api/auth/register} (public, no authentication required) accepts
a \texttt{RegisterRequest} DTO:
\begin{lstlisting}[language=Java, caption={RegisterRequest.java}]
@Data
public class RegisterRequest {
    private String username;
    private String password;
    private String role;   // DOCTOR | DIETICIAN | PATIENT | KITCHEN
}
\end{lstlisting}

Validation rules:
\begin{itemize}[noitemsep]
  \item \texttt{username} and \texttt{password} must not be blank.
  \item \texttt{role} must be one of \{DOCTOR, DIETICIAN, PATIENT, KITCHEN\};
        ADMIN is explicitly excluded from the allow-list.
  \item Duplicate usernames return \texttt{409 Conflict}.
\end{itemize}

On success the endpoint returns a \texttt{UserResponse} (\texttt{id},
\texttt{username}, \texttt{role}) with status \texttt{201 Created}.

\subsection{Frontend}
\texttt{RegisterModal.tsx} is a custom modal rendered from \texttt{LoginPage.tsx}
when the user clicks the \emph{Create Account} button.  It contains
username/password inputs and a role drop-down (ADMIN is hidden).
Server-side error messages (e.g.\ \emph{Username already taken}) are displayed
inline.  On success, \texttt{onSuccess(username, password)} is called, which
automatically submits the login form with the new credentials.

% ─────────────────────────────────────────────────────────────
\section{Feature: Recipe Images}

A \texttt{imageUrl} field (\texttt{String}) was added to the \texttt{Recipe}
entity and the \texttt{RecipeResponse} DTO.  \texttt{DataSeeder} populates all
15 seeded recipes with Unsplash CDN URLs (\texttt{?w=600\&q=80}).
The React \texttt{RecipeListPage} renders the URL as a full-width banner at the
top of each recipe card when the field is present.

% ─────────────────────────────────────────────────────────────
\section{Feature: Real-Time Notifications (WebSocket / STOMP)}

\subsection{Motivation}
Previously a patient could only order a recipe as a whole.  The updated system
lets the patient \emph{choose} which optional (modifiable) components to include,
while the system enforces that no selected component conflicts with any of the
patient's medical conditions.

\subsection{Backend Changes}

\subsubsection{\texttt{MealOrder} entity}
A new \texttt{ManyToMany} relationship stores the components the patient selected:

\begin{lstlisting}[language=Java, caption={MealOrder.java — selectedComponents field}]
@ManyToMany(fetch = FetchType.EAGER)
@JoinTable(
    name = "meal_order_selected_components",
    joinColumns = @JoinColumn(name = "meal_order_id"),
    inverseJoinColumns = @JoinColumn(name = "component_id")
)
@Builder.Default
private Set<RecipeComponent> selectedComponents = new HashSet<>();
\end{lstlisting}

\subsubsection{\texttt{MealOrderRequest} DTO}
\begin{lstlisting}[language=Java, caption={MealOrderRequest.java}]
@Data
public class MealOrderRequest {
    private Long recipeId;
    private List<Long> selectedComponentIds = Collections.emptyList();
}
\end{lstlisting}

\subsubsection{\texttt{MealOrderResponse} DTO}
A \texttt{List<ComponentResponse>} field named \texttt{selectedComponents} was
added so consumers can display exactly what was ordered.

\subsubsection{Validation in \texttt{MealOrderService.placeOrder()}}
The service validates in two passes:
\begin{enumerate}[noitemsep]
  \item \textbf{Main component}: if the recipe's non-modifiable main component
        conflicts with any of the patient's conditions, the order is rejected
        immediately with \texttt{400 Bad Request}.
  \item \textbf{Selected optional components}: for each ID in
        \texttt{selectedComponentIds} the service verifies:
        \begin{enumerate}[noitemsep]
          \item The component exists.
          \item It is a \texttt{ModifiableComponent}.
          \item It belongs to the requested recipe.
          \item None of its \texttt{incompatibleConditions} are in the patient's profile.
        \end{enumerate}
\end{enumerate}

\subsection{Frontend Changes}

\subsubsection{Type updates (\texttt{src/types/api.ts})}
\begin{lstlisting}[language=Java, caption={Updated TypeScript interfaces}]
export interface MealOrderRequest {
  recipeId: number;
  selectedComponentIds: number[];
}

export interface MealOrderResponse {
  // ... existing fields ...
  selectedComponents: ComponentResponse[];
}
\end{lstlisting}

\subsubsection{Inline component selector (\texttt{RecipeListPage.tsx})}
When a patient clicks \emph{Order Now} on a recipe that has at least one
modifiable component, an inline selector is displayed instead of immediately
submitting:
\begin{itemize}[noitemsep]
  \item Each modifiable component is shown as a labelled checkbox.
  \item Components whose \texttt{incompatibleConditionNames} intersect the
        patient's profile conditions are \emph{disabled} and annotated with the
        conflicting condition names.
  \item Compatible components can be freely toggled.
  \item \emph{Confirm Order} submits the order with the selected component IDs;
        \emph{Cancel} dismisses the selector without ordering.
  \item For recipes without modifiable components the order is submitted
        immediately as before.
\end{itemize}

\subsubsection{Order display}
Selected components are now shown in both \texttt{ActiveOrdersPage} and
\texttt{OrderHistoryPage} below each order card.

% ─────────────────────────────────────────────────────────────
\section{Feature: Meal Type Filter}

\subsection{Motivation}
Recipes are categorised as either a \emph{main} dish or a \emph{side} dish
(\texttt{MealType}: \texttt{MAIN} / \texttt{SIDE}).  There was previously no
dedicated UI control to filter by this dimension.

\subsection{Changes}

\subsubsection{Strategy layer}
A new \texttt{MealTypeSearchStrategy} class was added to
\texttt{searchStrategy.ts}:

\begin{lstlisting}[language=Java, caption={MealTypeSearchStrategy.ts}]
export class MealTypeSearchStrategy implements RecipeSearchStrategy {
  constructor(private readonly mealType: MealType) {}

  matches(recipe: RecipeResponse): boolean {
    return recipe.mealType === this.mealType;
  }
}
\end{lstlisting}

The \texttt{buildAndStrategy} factory was updated to accept an optional
\texttt{mealType} parameter, and pushes a \texttt{MealTypeSearchStrategy} into
the composite when a type is selected.

\subsubsection{UI}
A second row of tabs (All Types / Main / Side) was added directly below the
existing meal-course tabs in \texttt{RecipeListPage.tsx}.  Selecting a type
updates the \texttt{mealTypeFilter} state, which feeds into
\texttt{buildAndStrategy}, causing the recipe grid to update immediately
without a server round-trip.

% ─────────────────────────────────────────────────────────────
\section{Conflict Enforcement Summary}

The following table summarises where conflict enforcement takes place:

\begin{center}
\begin{tabular}{lll}
\toprule
\textbf{Enforcement point} & \textbf{Layer} & \textbf{Effect} \\
\midrule
Recipe card ribbon          & Frontend & Shows ``Conflicts'' warning on card \\
Order Now button            & Frontend & Disabled when main component conflicts \\
Component selector          & Frontend & Disables incompatible optional components \\
Main component check        & Backend  & 400 if main component conflicts patient conditions \\
Optional component check    & Backend  & 400 if any selected component conflicts \\
\bottomrule
\end{tabular}
\end{center}

% ─────────────────────────────────────────────────────────────
\section{API Endpoints}

\begin{center}
\begin{tabular}{lll}
\toprule
\textbf{Method} & \textbf{Path} & \textbf{Description} \\
\midrule
POST   & \texttt{/api/auth/register}          & Register new user (public, returns UserResponse) \\
GET    & \texttt{/api/auth/me}                & Current authenticated user info \\
POST   & \texttt{/api/orders}                 & Place order (Patient) \\
GET    & \texttt{/api/orders/active}          & Active orders for current patient \\
GET    & \texttt{/api/orders/history}         & Completed orders for current patient \\
GET    & \texttt{/api/recipes}                & List all recipes \\
GET    & \texttt{/api/recipes/\{id\}}          & Get single recipe \\
GET    & \texttt{/api/recipes/search}         & Search by name / component / conditions \\
POST   & \texttt{/api/recipes/\{id\}/check-compatibility} & Check against condition set \\
GET    & \texttt{/api/components}             & List all components \\
GET    & \texttt{/api/conditions}             & List all conditions \\
GET    & \texttt{/api/kitchen/orders}         & All active orders (Kitchen/Admin) \\
PUT    & \texttt{/api/kitchen/orders/\{id\}/advance} & Advance order status \\
GET    & \texttt{/api/doctor/patients}        & Doctor's patients (Doctor/Admin) \\
POST/DELETE & \texttt{/api/doctor/patients/\{pid\}/conditions/\{cid\}} & Assign/remove condition \\
GET    & \texttt{/api/notifications}          & Notifications for current user \\
PUT    & \texttt{/api/notifications/\{id\}/read} & Mark notification read \\
\bottomrule
\end{tabular}
\end{center}

% ─────────────────────────────────────────────────────────────
\section{Security}
All endpoints are secured with Spring Security HTTP Basic.
Role-based access control is enforced via \texttt{@PreAuthorize} annotations:
\begin{itemize}[noitemsep]
  \item Only \texttt{PATIENT} role may place or view orders.
  \item Only \texttt{DIETICIAN} / \texttt{ADMIN} may create, update, or delete
        recipes and components.
  \item \texttt{DOCTOR} / \texttt{ADMIN} may manage patient conditions and
        profiles.
  \item \texttt{KITCHEN} / \texttt{ADMIN} may advance order status.
  \item \texttt{/api/auth/**} is fully public (registration and login).
\end{itemize}
Passwords are hashed with BCrypt.  CSRF is disabled (stateless HTTP Basic API).
The \texttt{POST /api/auth/register} endpoint enforces a role allow-list
\{DOCTOR, DIETICIAN, PATIENT, KITCHEN\} — ADMIN accounts cannot be
self-registered.

% ─────────────────────────────────────────────────────────────
\section{Deployment}

\subsection{Docker Compose}
A \texttt{docker-compose.yml} at the project root orchestrates two services:
\begin{itemize}[noitemsep]
  \item \textbf{db} — \texttt{postgres:16-alpine}; data mounted on named volume
        \texttt{pgdata}; healthcheck via \texttt{pg\_isready}.
  \item \textbf{app} — multi-stage Docker image (Node 22 Alpine for the React
        build, Eclipse Temurin 21 JDK for the Maven build, JRE 21 runtime);
        receives \texttt{DATABASE\_URL / DATABASE\_USERNAME / DATABASE\_PASSWORD}
        as environment variables; \texttt{depends\_on: db: condition: service\_healthy}
        ensures the database is ready before the application starts.
\end{itemize}

To start the full stack:
\begin{lstlisting}[language=bash]
docker compose up --build
\end{lstlisting}
The application is available at \texttt{http://localhost:8081} and
Swagger UI at \texttt{http://localhost:8081/swagger-ui.html}.

\subsection{Local Development}
For iterative development without Docker, a local PostgreSQL instance is required:
\begin{lstlisting}[language=bash]
docker run -d --name pg \
  -e POSTGRES_DB=recipemaker \
  -e POSTGRES_USER=recipemaker \
  -e POSTGRES_PASSWORD=recipemaker \
  -p 5432:5432 postgres:16-alpine

cd backend && ./mvnw spring-boot:run
cd frontend && npm install && npm run dev
\end{lstlisting}
The Vite dev server proxies \texttt{/api} and \texttt{/ws} to
\texttt{localhost:8081}.

\subsection{Running Tests}
\begin{lstlisting}[language=bash]
cd backend && ./mvnw test   # 120 tests, H2 in-memory, no DB required
\end{lstlisting}

% ─────────────────────────────────────────────────────────────
\section{Conclusion}
The Recipe Maker application provides a complete, role-aware hospital meal
ordering workflow.  Patients browse recipe cards (with images), select optional
components in a compatibility-aware modal, and receive real-time status updates
via WebSocket push notifications.  Kitchen staff advance order statuses from a
dedicated queue view.  Doctors manage patient conditions that instantly affect
ordering eligibility.  Dieticians maintain the recipe and component catalogue.

All business rules are enforced at both layers: the frontend disables
incompatible options immediately based on the loaded patient profile, and the
backend re-validates every selection before persisting the order.  The search
and filter system uses the Strategy pattern on both tiers, keeping the codebase
extensible — adding a new filter criterion requires only a new leaf-strategy
class and a one-line addition to \texttt{buildAndStrategy}.  The Observer
pattern on the backend similarly allows new notification channels (email, SMS)
to be wired in without touching the core order-management service.

\end{document}
